\answer{ex:21:1}{(а)~$T=\tau_{\text{вос}}\ln\frac{1-x_0}{1-x^*}$. (б)~$1.84$~с и $3.68$~с; $\dd T/\dd x^*=\tau_{\text{вос}}/(1-x^*)=8$~с и $80$~с. (в)~От $3.36$ до $4.24$~с: у $x^*\to1$ залпы нерегулярны.}
\answer{ex:21:2}{(а)~$10^{-4}$~Вт; для мозга $8.6$~Вт, как в~\eqref{eq:energy}. (б)~$205$~Мбит/с, $0.2$~Вт --- в $2\cdot10^3$ раз больше культуры.}
\answer{ex:21:3}{(а)~$x_{\text{уст}}=1/(1+Ur_b\tau_{\text{вос}})$; $r_b>(1/x^*-1)/(U\tau_{\text{вос}})=0.28$~Гц при $x^*=0.9$ и $0.025$~Гц при $x^*=0.99$. (б)~Сила синапса падает до $x_{\text{уст}}$, то есть на $1$--$10\%$. (в)~Залп запускается там, где запас восстановлен.}
\answer{ex:21:4}{$u(t-k)$ ортонормированы, $m_k=\|P_Vu(t-k)\|^2$, где $P_V$ --- проектор на $N$-мерную оболочку выходов; по неравенству Бесселя $\sum_k m_k\le N$.}
\answer{ex:21:5}{(а)~$\mathrm{MC}\approx9$--$11$ с $\tanh$, $15$--$18$ и $17$--$21$ для линейного при $0.9$ и $0.99$ (три затравки); всегда $\le30$. (б)~Резервуар: точность $0.985$--$1.0$, $r^2=0.84$--$0.91$; отводы и резервуар без смещений: $r^2<0.01$ --- нечётная нелинейность не создаёт чётных произведений.}
\answer{ex:21:6}{(а)~При $P\le n$ разделимы все $2^P$ разметок. (б)~$n\le102$. (в)~На $5.6$ стандартных отклонения выше $50\%$.}
\answer{ex:21:7}{Открытая задача. Ключевой контроль: заморозить считывание и проверить, сохраняется ли улучшение после паузы без стимуляции; если нет, изменилось состояние, а не веса.}
